% GENERATED FILE. Edit canonical lessons, then run npm run generate:books.
% canonical-source-hash: fnv1a64:385211d3d3b0c616
% canonical-lessons: KA-C243-gurutisu, KA-C243-uccarisu, KA-C243-kudisu, KA-C243-nibhayisu, KA-C243-alladisu, KA-R243-first-pass-words-from-chapters-235-238, KA-R243-second-pass-words-from-chapters-239-243

\chapter{To recognise, To pronounce, To add up, To manage, To shake}
\label{ch:ka-to-recognise-to-pronounce-to-add-up-to-manage-to-shake}
\hlchaptermodality{243}

\begin{chapteropening}
\textbf{By the end of this chapter:} \emph{I can say to recognise, to pronounce, to add up, to manage and to shake.}

Complete the last lesson of chapter 243: I can say to recognise, to pronounce, to add up, to manage and to shake.
\end{chapteropening}

\section[\texorpdfstring{\kn{ಗುರುತಿಸು} (gurutisu)}{ಗುರುತಿಸು (gurutisu)}]{\kn{ಗುರುತಿಸು} (gurutisu) — to recognise}

\label{lesson:KA-C243-gurutisu}



\begin{quote}
Before the new one: say the Kannada for to dip, then the Kannada for to tear.
\end{quote}

\subsection*{You'll want to know: \kn{ಗುರುತಿಸು}}
\textbf{\kn{ಗುರುತಿಸು}} — \emph{gurutisu} — \textquotedblleft{}to recognise\textquotedblright{}.

\begin{grammarlens}[title={What you've built}]
One more word for this chapter.
\end{grammarlens}

\subsection*{Your turn}
\begin{itemize}
\raggedright
  \item \emph{Say it:} \emph{gurutisu}
  \item \emph{Say it:} \emph{gurutisu}, once more
  \item \emph{Say it:} say \emph{gurutisu}
  \item \emph{Recall:} say the Kannada for to dip, then the Kannada for to tear, then say \emph{gurutisu} again
\end{itemize}

\begin{tcolorbox}[breakable,colback=teal!4,colframe=teal!35!black,title={Before you move on}]
What does \kn{ಗುರುತಿಸು} mean? (\textquotedblleft{}To recognise\textquotedblright{}.) Say it once more.
\end{tcolorbox}
\section[\texorpdfstring{\kn{ಉಚ್ಚರಿಸು} (uccarisu)}{ಉಚ್ಚರಿಸು (uccarisu)}]{\kn{ಉಚ್ಚರಿಸು} (uccarisu) — to pronounce}

\label{lesson:KA-C243-uccarisu}



\begin{quote}
Before the new one: say the Kannada for to tear, then the Kannada for to recognise.
\end{quote}

\subsection*{You'll want to know: \kn{ಉಚ್ಚರಿಸು}}
\textbf{\kn{ಉಚ್ಚರಿಸು}} — \emph{uccarisu} — \textquotedblleft{}to pronounce\textquotedblright{}.

\begin{grammarlens}[title={What you've built}]
One more word for this chapter.
\end{grammarlens}

\subsection*{Your turn}
\begin{itemize}
\raggedright
  \item \emph{Say it:} \emph{uccarisu}
  \item \emph{Say it:} \emph{uccarisu}, once more
  \item \emph{Say it:} say \emph{uccarisu}
  \item \emph{Recall:} say the Kannada for to tear, then the Kannada for to recognise, then say \emph{uccarisu} again
\end{itemize}

\begin{tcolorbox}[breakable,colback=teal!4,colframe=teal!35!black,title={Before you move on}]
What does \kn{ಉಚ್ಚರಿಸು} mean? (\textquotedblleft{}To pronounce\textquotedblright{}.) Say it once more.
\end{tcolorbox}
\section[\texorpdfstring{\kn{ಕೂಡಿಸು} (kūḍisu)}{ಕೂಡಿಸು (kūḍisu)}]{\kn{ಕೂಡಿಸು} (kūḍisu) — to add up (numbers)}

\label{lesson:KA-C243-kudisu}



\begin{quote}
Before the new one: say the Kannada for to recognise, then the Kannada for to pronounce.
\end{quote}

\subsection*{You'll want to know: \kn{ಕೂಡಿಸು}}
\textbf{\kn{ಕೂಡಿಸು}} — \emph{kūḍisu} — \textquotedblleft{}to add up (numbers)\textquotedblright{}.

\begin{grammarlens}[title={What you've built}]
One more word for this chapter.
\end{grammarlens}

\subsection*{Your turn}
\begin{itemize}
\raggedright
  \item \emph{Say it:} \emph{kūḍisu}
  \item \emph{Say it:} \emph{kūḍisu}, once more
  \item \emph{Say it:} say \emph{kūḍisu}
  \item \emph{Recall:} say the Kannada for to recognise, then the Kannada for to pronounce, then say \emph{kūḍisu} again
\end{itemize}

\begin{tcolorbox}[breakable,colback=teal!4,colframe=teal!35!black,title={Before you move on}]
What does \kn{ಕೂಡಿಸು} mean? (\textquotedblleft{}To add up (numbers)\textquotedblright{}.) Say it once more.
\end{tcolorbox}
\section[\texorpdfstring{\kn{ನಿಭಾಯಿಸು} (nibhāyisu)}{ನಿಭಾಯಿಸು (nibhāyisu)}]{\kn{ನಿಭಾಯಿಸು} (nibhāyisu) — to manage, to cope with}

\label{lesson:KA-C243-nibhayisu}



\begin{quote}
Before the new one: say the Kannada for to pronounce, then the Kannada for to add up.
\end{quote}

\subsection*{You'll want to know: \kn{ನಿಭಾಯಿಸು}}
\textbf{\kn{ನಿಭಾಯಿಸು}} — \emph{nibhāyisu} — \textquotedblleft{}to manage, to cope with\textquotedblright{}.

\begin{grammarlens}[title={What you've built}]
One more word for this chapter.
\end{grammarlens}

\subsection*{Your turn}
\begin{itemize}
\raggedright
  \item \emph{Say it:} \emph{nibhāyisu}
  \item \emph{Say it:} \emph{nibhāyisu}, once more
  \item \emph{Say it:} say \emph{nibhāyisu}
  \item \emph{Recall:} say the Kannada for to pronounce, then the Kannada for to add up, then say \emph{nibhāyisu} again
\end{itemize}

\begin{tcolorbox}[breakable,colback=teal!4,colframe=teal!35!black,title={Before you move on}]
What does \kn{ನಿಭಾಯಿಸು} mean? (\textquotedblleft{}To manage, to cope with\textquotedblright{}.) Say it once more.
\end{tcolorbox}
\section[\texorpdfstring{\kn{ಅಲ್ಲಾಡಿಸು} (allāḍisu)}{ಅಲ್ಲಾಡಿಸು (allāḍisu)}]{\kn{ಅಲ್ಲಾಡಿಸು} (allāḍisu) — to shake (something)}

\label{lesson:KA-C243-alladisu}



\begin{quote}
Before the new one: say the Kannada for to add up, then the Kannada for to manage.
\end{quote}

\subsection*{You'll want to know: \kn{ಅಲ್ಲಾಡಿಸು}}
\textbf{\kn{ಅಲ್ಲಾಡಿಸು}} — \emph{allāḍisu} — \textquotedblleft{}to shake (something)\textquotedblright{}.

\begin{grammarlens}[title={What you've built}]
One more word for this chapter.
\end{grammarlens}

\subsection*{Your turn}
\begin{itemize}
\raggedright
  \item \emph{Say it:} \emph{allāḍisu}
  \item \emph{Say it:} \emph{allāḍisu}, once more
  \item \emph{Say it:} say \emph{allāḍisu}
  \item \emph{Recall:} say the Kannada for to add up, then the Kannada for to manage, then say \emph{allāḍisu} again
\end{itemize}

\begin{tcolorbox}[breakable,colback=teal!4,colframe=teal!35!black,title={Before you move on}]
What does \kn{ಅಲ್ಲಾಡಿಸು} mean? (\textquotedblleft{}To shake (something)\textquotedblright{}.) Say it once more.
\end{tcolorbox}
\section[(dialogue)]{Words from chapters 235-238 — a first pass}

\label{lesson:KA-R243-first-pass-words-from-chapters-235-238}



\begin{quote}
Say the words for to manage and to shake.
\end{quote}

\subsection*{Your turn}
Say these aloud:

\begin{itemize}
\raggedright
  \item to fall, to return, to touch, to lift, to throw
  \item to untie, to wake, to wake up, to bargain, to climb
  \item to bite, to use, to repair, to hit, to try
  \item to type, to get lost, to agree, to refuse, to earn
\end{itemize}

\begin{tcolorbox}[breakable,colback=teal!4,colframe=teal!35!black,title={Before you move on}]
To fall? (\textbf{\kn{ಬೀಳು}}, \emph{bīḷu}.) To return? (\textbf{\kn{ಹಿಂದಿರುಗು}}, \emph{hindirugu}.) To touch? (\textbf{\kn{ಮುಟ್ಟು}}, \emph{muṭṭu}.) To lift? (\textbf{\kn{ಎತ್ತು}}, \emph{ettu}.) To throw? (\textbf{\kn{ಎಸೆ}}, \emph{ese}.) To untie? (\textbf{\kn{ಬಿಚ್ಚು}}, \emph{biccu}.) To wake? (\textbf{\kn{ಎಬ್ಬಿಸು}}, \emph{ebbisu}.) To wake up? (\textbf{\kn{ಎಚ್ಚರಗೊಳ್ಳು}}, \emph{eccaragoḷḷu}.) To bargain? (\textbf{\kn{ಚೌಕಾಸಿ} \kn{ಮಾಡು}}, \emph{caukāsi māḍu}.) To climb? (\textbf{\kn{ಏರು}}, \emph{ēru}.) To bite? (\textbf{\kn{ಕಚ್ಚು}}, \emph{kaccu}.) To use? (\textbf{\kn{ಬಳಸು}}, \emph{baḷasu}.) To repair? (\textbf{\kn{ಸರಿಪಡಿಸು}}, \emph{saripaḍisu}.) To hit? (\textbf{\kn{ಹೊಡೆ}}, \emph{hoḍe}.) To try? (\textbf{\kn{ಪ್ರಯತ್ನಿಸು}}, \emph{prayatnisu}.) To type? (\textbf{\kn{ಟೈಪ್} \kn{ಮಾಡು}}, \emph{ṭaip māḍu}.) To get lost? (\textbf{\kn{ಕಳೆದುಹೋಗು}}, \emph{kaḷeduhōgu}.) To agree? (\textbf{\kn{ಒಪ್ಪು}}, \emph{oppu}.) To refuse? (\textbf{\kn{ನಿರಾಕರಿಸು}}, \emph{nirākarisu}.) To earn? (\textbf{\kn{ಸಂಪಾದಿಸು}}, \emph{sampādisu}.)
\end{tcolorbox}
\section[(dialogue)]{Words from chapters 239-243 — a second pass}

\label{lesson:KA-R243-second-pass-words-from-chapters-239-243}



\begin{quote}
Say the words for to celebrate and to decorate.
\end{quote}

\subsection*{Your turn}
Say these aloud:

\begin{itemize}
\raggedright
  \item to celebrate, to decorate, to share out, to pay a visit, to hug
  \item to bark, to shout, to sink, to argue, to imagine
  \item to criticise, to share, to hate, to find out, to borrow
  \item to temper, to soak, to wring, to dip, to tear
  \item to recognise, to pronounce, to add up, to manage, to shake
\end{itemize}

\begin{tcolorbox}[breakable,colback=teal!4,colframe=teal!35!black,title={Before you move on}]
To celebrate? (\textbf{\kn{ಆಚರಿಸು}}, \emph{ācarisu}.) To decorate? (\textbf{\kn{ಅಲಂಕರಿಸು}}, \emph{alaṅkarisu}.) To share out? (\textbf{\kn{ಹಂಚು}}, \emph{hañcu}.) To pay a visit? (\textbf{\kn{ಭೇಟಿ} \kn{ಕೊಡು}}, \emph{bhēṭi koḍu}.) To hug? (\textbf{\kn{ಅಪ್ಪಿಕೊ}}, \emph{appiko}.) To bark? (\textbf{\kn{ಬೊಗಳು}}, \emph{bogaḷu}.) To shout? (\textbf{\kn{ಕೂಗು}}, \emph{kūgu}.) To sink? (\textbf{\kn{ಮುಳುಗು}}, \emph{muḷugu}.) To argue? (\textbf{\kn{ವಾದಿಸು}}, \emph{vādisu}.) To imagine? (\textbf{\kn{ಕಲ್ಪಿಸಿಕೊ}}, \emph{kalpisiko}.) To criticise? (\textbf{\kn{ಟೀಕಿಸು}}, \emph{ṭīkisu}.) To share? (\textbf{\kn{ಹಂಚಿಕೊಳ್ಳು}}, \emph{hañcikoḷḷu}.) To hate? (\textbf{\kn{ದ್ವೇಷಿಸು}}, \emph{dvēṣisu}.) To find out? (\textbf{\kn{ಕಂಡುಹಿಡಿ}}, \emph{kaṇḍuhiḍi}.) To borrow? (\textbf{\kn{ಸಾಲ} \kn{ಪಡೆ}}, \emph{sāla paḍe}.) To temper? (\textbf{\kn{ಒಗ್ಗರಣೆ} \kn{ಹಾಕು}}, \emph{oggaraṇe hāku}.) To soak? (\textbf{\kn{ನೆನೆಸು}}, \emph{nenesu}.) To wring? (\textbf{\kn{ಹಿಂಡು}}, \emph{hiṇḍu}.) To dip? (\textbf{\kn{ಅದ್ದು}}, \emph{addu}.) To tear? (\textbf{\kn{ಹರಿ}}, \emph{hari}.) To recognise? (\textbf{\kn{ಗುರುತಿಸು}}, \emph{gurutisu}.) To pronounce? (\textbf{\kn{ಉಚ್ಚರಿಸು}}, \emph{uccarisu}.) To add up? (\textbf{\kn{ಕೂಡಿಸು}}, \emph{kūḍisu}.) To manage? (\textbf{\kn{ನಿಭಾಯಿಸು}}, \emph{nibhāyisu}.) To shake? (\textbf{\kn{ಅಲ್ಲಾಡಿಸು}}, \emph{allāḍisu}.)
\end{tcolorbox}
